\scanpage{14}
\mspage{4}
\deleted{\unsure{Pic}} additive ordinaire formée des triples $(M,N,\sigma\colon M\longrightarrow{}_2N)$. Si $\unsure{\mathscr P|P},\mathscr P'$ sont associés à $(M,N,\sigma),(M',N',\sigma')$, alors deux foncteurs $\Lambda,\Lambda'\colon\mathscr P\longrightarrow\mathscr P'$ sont isom. ssi. ils définissent le $=$ hom $(M,N,\sigma)\longrightarrow(M',N',\sigma)$ (donc $D$ ``est'' la catégorie \unsure{déduite|discrète} de $(\mathrm{Pic})$ en passant au quotient par les 2-flèches).
% slanted
\margin{NB $\underline{\operatorname{Hom}}_{\otimes\,\mathrm{AUC}}(\mathscr P,\mathscr P')$ est un groupoïde}
Mais bien sûr, la catégorie des $\otimes$-foncteur AUC $\mathscr P\longrightarrow\mathscr P'$ qui sont subordonnés à un hom donné $(M,N,\sigma)\mathbin{\unsure{-|\to}}(M',N',\sigma)$ forment un groupoïde qui est un torseur sous la Pic-catégorie définie par \unsure{tronsection|transection|troncature|troncation} à partir de $R\underline{\operatorname{Hom}}(M,N')$ (d'invariant $\operatorname{Ext}^{\unsure{1}}(M,N'),\operatorname{Hom}(M,N')$.

7) \deleted{\unsure{Présentation|Il résulte|Traductions}} \uline{Extension aux champs de Picard sur un topos}

Ils forment une 2-\inserted{$\otimes$-}catégorie AUC, et \inserted{définissent} un 2-$\unsure{\otimes|O}$-Champ AUC. Il s'envoie dans la 1-\inserted{$\otimes$-}cat. AUC (ou 1-$\otimes$-champ AUC) des $(M,N,\sigma\colon M\longrightarrow{}_2N)$, mais si $\Lambda,\Lambda'\colon\mathscr P\longrightarrow\mathscr P'$ définissent le $=$
\[
(M,N,\sigma)\xrightarrow{\lambda}(M',N',\sigma'),
\]
\deleted{\unsure{on ne peut plus}}\inserted{(p. ex. l'identité)}, on
% archive sheet number: 13
